\question[10]
Available evidence suggests that the last five major mass extinctions have been caused by different events (e.g., asteroids, volcanism) but the extinctions ultimately share a common cause. What is that cause? Describe one type of evidence (from the text or from the Barnosky et al. paper) that suggests the currently changing climate might be close to causing a sixth mass extinction.
	\begin{solution}
	Answer.
	\end{solution}	
\vspace*{\stretch{1}}


\question
\fullwidth{Male swordtail characin fishes (\textit{Corynopoma riisei}) have a long, narrow extension of the gill cover. The extension has a ``flag'' on the tip, which mimics the fish's prey. The male uses the flag to lure females to the male for mating. A study by Arnqvist and Kolm (2010) showed that the shape of the flag was strongly correlated with environmental variables like prey availability. Based on this information, answer the next two questions.}
	\begin{parts}
	\part[5]
	Explain whether this is intrasexual or intersexual selection and why.
	
	\AnswerBox{1\baselineskip}{%
	Intersexual. Females are choosing males based on a characteristic of the male.
	}
	
	\part[10]
	Decide whether the flag more likely evolved by sensory bias or sensory drive. Explain what is sensory bias or sensory drive (depending on your choice), and justify your decision by explaining how your choice is consistent with the scenario described above.

	\AnswerBox{1\baselineskip}{%
	Sensory drive. Sensory drive is when female drives evolution of a trait based on the ability of the trait to stand out from the environment. Females are choosing males whose trait indicates something about the environment.
	}
	\end{parts}

\question[5]
Explain Muller's ratchet. Be sure you covered in lecture.

\quesiton[5]
Explain the Red Queen hypothesis as it relates to the evolution of sexual reproduction.

\question[10]
Male satin bowerbirds decorate their bowers with bright blue objects they find around the area. The bower is not a nest. It's a structure built solely to attract females. Females build the nest elsewhere. Assuming the male does \emph{not} help raise his offspring, explain how the male's ability to gather decorative objects may serve as an indirect indicator of male genetic quality that affects female choice. In your explanation, be sure to state which intersexual selection hypothesis best fits this scenario, and why.
	\begin{solution}
	Good genes hypothesis. Males with the most decorated bowers expended more energy to find the objects. That would also make them more vulnerable to predation as they spend more time searching. The most decorated bowers would indicate to the female that the male has good genetic quality because he is still alive.
	\end{solution}
	

\question[20]% Make 15 point question for future. 
Discuss in terms of fitness and adaptive flexibility the advantages and disadvantages of both asexual and sexual reproduction. That is, what are the advantages and disadvantages of maintaining the same genome (asexual) versus a variable genome (sexual) in relation to the current environment (in a broad sense) and possible future changes to the environment. Be sure to include and explain the ``cost of sex'' in your discussion, as it relates to advantages of disadvantages of each mode of reproduction.
%NEW VERSION This could be 20 point question, but make the cost of sex advantage worth 5 points.
% Explain the ``cost of sex'' as it relates to the fitness advantage or disadvantage of asexual and sexual reproduction. Then explain one additional advantage and disadvantage of maintaining the same genome (asexual reproduction), and one advantage and disadvantage of maintaining a variable genome (sexual reproduction) in relation to the current environment and possible future changes to the environment. Think of environment in a broad sense, including parasites/disease.
	\begin{solution}
	3 pts: advantages of asexual.
	
	3 pts: disadvantages of asexal. \vspace{\baselineskip}
	
	3 pts: advantages of sexual.
	
	3 pts: disadvantages of sexual.
	
	3 pts: cost of sex.
	\end{solution}
